% Part: applied-modal-logics
% Chapter: epistemic-logic
% Section: bisimulations

\documentclass[../../../include/open-logic-section]{subfiles}

\begin{document}

\olfileid{aml}{el}{bsd}

\olsection{دوشبیه‌سازی‌ها}

یک پرسش باقی‌مانده دربارهٔ توان بیانی زبان معرفتی ما به رابطهٔ میان
مدل‌ها و فرمول‌هایی مربوط می‌شود که در آن‌ها برقرارند. از نتایج تناظر
قاب دیدیم که وقتی فرمول‌های معینی در یک قاب معتبر باشند، تضمین می‌کنند
که آن قاب‌ها خواص معینی را برآورده سازند. اما آیا زبان موجهاتی ما، برای
نمونه، اجازه می‌دهد میان جهانی با یک پیکان بازتابی و زنجیره‌ای نامتناهی
از جهان‌ها که هریک به جهان بعدی می‌انجامد تمایز بگذاریم؟ یعنی آیا
فرمول~$A$ای وجود دارد که تنها در یکی از این دو جهان برقرار باشد؟

دوشبیه‌سازی رابطه‌ای است که می‌توانیم میان مدل‌های رابطه‌ای تعریف کنیم
تا بگوییم ساختار آن‌ها عملاً یکسان است. چنان‌که خواهیم دید، این رابطه
معنایی از هم‌ارزی میان مدل‌ها را به دست می‌دهد که در زبان معرفتی ما
قابل بیان است.

\begin{defn}[دوشبیه‌سازی]
فرض کنید $M_1 = \tuple{W_1, R_1, V_1}$ و
$M_2 = \tuple{W_2, R_2, V_2}$ دو مدل رابطه‌ای باشند و
$\mathcal{R} \subseteq W_1 \times W_2$ رابطه‌ای دوتایی باشد. می‌گوییم
$\mathcal{R}$ یک \emph{دوشبیه‌سازی} است هرگاه برای هر
$\tuple{w_1, w_2} \in \mathcal{R}$، شرایط زیر برقرار باشند:

\begin{enumerate}
  \item برای هر !!{propositional variable}~$p$،
  $w_1 \in V_1(p)$ اگر و تنها اگر~$w_2 \in V_2(p)$.
  \item برای همهٔ عامل‌های~$a \in G$ و جهان‌های~$v_1 \in W_1$، اگر
  $R_{1,a} w_1 v_1$، آنگاه~$v_2 \in W_2$ای وجود دارد چنان‌که
  $R_{2,a} w_2 v_2$ و~$\tuple{v_1, v_2} \in \mathcal{R}$.
  \item برای همهٔ عامل‌های~$a \in G$ و جهان‌های~$v_2 \in W_2$، اگر
   $R_{2,a} w_2 v_2$، آنگاه~$v_1 \in W_1$ای وجود دارد چنان‌که
   $R_{1,a} w_1 v_1$ و~$\tuple{v_1, v_2} \in \mathcal{R}$.
\end{enumerate}

هرگاه دوشبیه‌سازی‌ای میان~$M_1$ و~$M_2$ وجود داشته باشد که جهان‌های
$w_1$ و~$w_2$ را به هم مرتبط کند، می‌توانیم بنویسیم
$\tuple{M_1, w_1} \leftrightarroweq \tuple{M_2, w_2}$ و
$\tuple{M_1, w_1}$ و~$\tuple{M_2, w_2}$ را \emph{دوشبیه} بنامیم.
\end{defn}

بندهای مختلف رابطهٔ دوشبیه‌سازی امور متفاوتی را تضمین می‌کنند. بند~1
تضمین می‌کند جهان‌های دوشبیه فرمول‌های فاقد وجه یکسانی را ارضا کنند،
زیرا توافق دربارهٔ همهٔ !!{propositional variable}ها را تضمین می‌کند.
دو بند دیگر، که به‌ترتیب گاه بندهای «پیش» و «پس» نامیده می‌شوند، تضمین
می‌کنند روابط دسترسی‌پذیری ساختار یکسانی داشته باشند.

\begin{thm}
اگر~$\tuple{M_1, w_1} \leftrightarroweq \tuple{M_2, w_2}$، آنگاه برای
هر !!{formula}~$!A$، $\mSat{M_1}{!A}[w_1]$ اگر و تنها اگر
$\mSat{M_2}{!A}[w_2]$.
\end{thm}

\begin{figure}
  \begin{center}
    \begin{tikzpicture}[modal]
      \node[world] (w1) {$w_1$};
      \node[world] (w2) [above left=of w1]{$w_2$};
      \node[world] (w3) [above right=of w1] {$w_3$};
      \draw[<->] (w1) to node {$a$} (w2);
      \draw[->,loop below] (w1) to node {$a$} (w1);
      \draw[<->] (w1) to [swap] node {$a$} (w3);
      \draw[->,loop above] (w2) to node {$a$} (w2) ;
      \draw[->,loop above] (w3) to node {$a$} (w3) ;
      
      \node[world](v1)[right of=w1, xshift=1.5in]{$v_1$};
      \node[world](v2)[above of=v1]{$v_2$};
      \draw[<->] (v1) to node {$a$} (v2);
      \draw[->,loop below] (v1) to node {$a$} (v1);
      \draw[->,loop above] (v2) to node {$a$} (v2) ;

      \draw[-,dotted] (w1) to (v1) ;
      \draw[-,dotted,bend left=45] (w2) to (v2) ;
      \draw[-,dotted,bend left=30] (w3) to (v2) ;
    \end{tikzpicture}
  \end{center}
  \caption{دو مدل دوشبیه.}
  \ollabel{fig:bisimilar}
\end{figure}

هرچند دو مدلِ تصویرشده در~\olref{fig:bisimilar} کاملاً یکسان نیستند،
دوشبیه‌سازی‌ای وجود دارد که جهان‌های~$w_1$ و~$v_1$ را به هم مرتبط
می‌کند. این دوشبیه‌سازی هم~$w_2$ و هم~$w_3$ را نیز به~$v_2$ مرتبط
می‌کند؛ زیرا هیچ چیزِ بیان‌پذیری در زبان موجهاتی ما نیست که واقعاً
بتواند میان آن‌ها تمایز بگذارد. البته اگر~$w_2$ و~$w_3$،
!!{propositional variable}های متفاوتی را ارضا می‌کردند، وضع متفاوت بود.

\end{document}
